\documentclass[11pt,letterpaper]{article}

\usepackage[utf8]{inputenc}
\usepackage[T1]{fontenc}
\usepackage[spanish,es-noshorthands]{babel}
\usepackage{amsmath,amssymb,amsthm}
\usepackage{enumitem}
\usepackage{booktabs}
\usepackage[margin=2.5cm]{geometry}
\usepackage[hidelinks]{hyperref}

\theoremstyle{definition}
\newtheorem{definicion}{Definición}
\newtheorem{restriccion}{Restricción}
\theoremstyle{remark}
\newtheorem*{obs}{Observación}

\newcommand{\N}{\mathbb{N}}
\newcommand{\R}{\mathbb{R}}
\newcommand{\card}[1]{\lvert #1 \rvert}
\newcommand{\pend}[1]{\textbf{[POR DEFINIR: #1]}}

\title{Algoritmo de asignación de horarios estudiantiles\\
\large Definición formal del problema --- Entrega Semana 10}
\author{Análisis de Algoritmos --- Proyecto 2026-30\\
Pontificia Universidad Javeriana\\[2pt]
\textit{Integrantes: Gustavo Aguilar, Santiago Hernández}}
\date{}

\begin{document}
\maketitle

%--------------------------------------------------------------------
\section{Alcance y supuestos}
%--------------------------------------------------------------------
Este documento formaliza el problema de asignar, a cada estudiante de un
conjunto dado, un horario semanal de clases. La formalización es
\emph{independiente del algoritmo}: sólo define qué es una instancia
(entrada), qué es una solución (salida), cuándo una solución es
\emph{factible} y qué significa que sea \emph{mejor} que otra.

Supuestos de modelado:
\begin{enumerate}[label=(S\arabic*)]
  \item El tiempo se discretiza en franjas de duración fija; el horario es
        semanal y se repite igual todas las semanas del semestre.
  \item La oferta académica (grupos, profesores, salones, franjas y cupos)
        está \emph{fija} y es dato de entrada; el algoritmo no la modifica.
  \item Cada estudiante declara qué materias puede cursar (ya incluye
        prerrequisitos aprobados); el algoritmo no valida el plan de estudios.
  \item Las preferencias subjetivas (profesores, amigos, tiempos libres)
        llegan ya cuantificadas como parámetros numéricos de la instancia.
  \item El horario de cada estudiante se evalúa por separado, salvo en el
        criterio social y en los cupos, que acoplan a los estudiantes.
\end{enumerate}

%--------------------------------------------------------------------
\section{Entrada: definición de instancia}
%--------------------------------------------------------------------

\subsection{Conjuntos base}
\begin{definicion}[Conjuntos]
\begin{itemize}[nosep]
  \item $E$: estudiantes. \quad $C$: materias. \quad $G$: grupos (secciones).
  \item $P$: profesores. \quad $R$: salones.
  \item $D=\{1,\dots,n_d\}$: días de la semana (típicamente $n_d=6$).
  \item $H=\{1,\dots,n_h\}$: periodos dentro de un día, de duración $L$ minutos.
  \item $S = D\times H$: franjas semanales. Para $s=(\delta,h)$ se escribe
        $\mathrm{dia}(s)=\delta$ y $\mathrm{per}(s)=h$.
\end{itemize}
Todos los conjuntos son finitos.
\end{definicion}

\subsection{Oferta académica}
Cada grupo $g\in G$ tiene:
\begin{itemize}[nosep]
  \item $c(g)\in C$: materia a la que pertenece. Se define $G_c=\{g\in G : c(g)=c\}$.
  \item $p(g)\in P$: profesor que lo dicta.
  \item $q(g)\in\N$: cupo máximo.
  \item $\Sigma(g)\subseteq S\times R$: \emph{sesiones}. Un par $(s,r)\in\Sigma(g)$
        indica que el grupo ocupa la franja $s$ en el salón $r$. Cada franja
        aparece a lo sumo una vez en $\Sigma(g)$. Se define
        $T(g)=\{s : (s,r)\in\Sigma(g)\}$ y, para $s\in T(g)$, $\mathrm{sal}_g(s)$
        el salón correspondiente.
\end{itemize}
Cada materia $c\in C$ tiene $\mathrm{cr}(c)\in\N$ créditos.

\subsection{Logística}
Función de distancia entre salones $d:R\times R\to\N$ (minutos de
desplazamiento a pie), con $d(r,r)=0$ y $d(r,r')=d(r',r)$. Parámetro
$\beta\in\N$: minutos de descanso reglamentarios entre periodos consecutivos.

\subsection{Datos por estudiante}
Para cada $e\in E$:
\begin{itemize}[nosep]
  \item $A(e)\subseteq C$: materias que puede cursar (elegibles).
  \item $\rho(e,c)\in\{1,\dots,\rho_{\max}\}$ para $c\in A(e)$: \emph{prioridad}
        de cursar $c$ (mayor = más necesaria; una materia obligatoria tiene
        prioridad alta).
  \item $U(e)\subseteq S$: franjas \emph{no disponibles} (trabajo, etc.).
  \item $F(e)\subseteq S$: franjas que \emph{prefiere dejar libres}.
  \item $\kappa_{\min}(e)\le\kappa_{\max}(e)$: rango de créditos admisible.
  \item $\pi(e,p)\in[-1,1]$ para $p\in P$: empatía con el profesor $p$
        ($-1$ rechazo, $0$ neutro, $1$ afinidad).
\end{itemize}
Afinidad social entre estudiantes: $\omega:E\times E\to[0,1]$, simétrica y
con $\omega(e,e)=0$. Un valor alto indica que $e$ y $e'$ desean compartir
clase (amigos o grupo de trabajo).

\subsection{Parámetros de evaluación}
Pesos $\lambda_1,\dots,\lambda_5\ge 0$ con $\sum_i\lambda_i=1$ (importancia
relativa de los criterios) y constantes de normalización $N_1,\dots,N_5>0$
(ver §\ref{sec:objetivo}).

\begin{definicion}[Instancia]
Una instancia es la tupla
\[
\mathcal{I}=\bigl(E,C,G,P,R,D,H,L,\ c,p,q,\Sigma,\mathrm{cr},\ d,\beta,\
A,\rho,U,F,\kappa_{\min},\kappa_{\max},\pi,\omega,\ \lambda,N\bigr).
\]
El \emph{tamaño} de la instancia se mide por $\card{E}$, $\card{G}$ y
$\sum_{g}\card{\Sigma(g)}$.
\end{definicion}

%--------------------------------------------------------------------
\section{Salida: definición de solución}
%--------------------------------------------------------------------
\begin{definicion}[Asignación]
Una asignación es un conjunto $X\subseteq E\times G$. Se escribe
$X_e=\{g : (e,g)\in X\}$ (horario del estudiante $e$) y
$X^g=\{e : (e,g)\in X\}$ (inscritos en $g$). Equivalentemente, variables
binarias $x_{e,g}=1 \iff (e,g)\in X$.
\end{definicion}

La salida del algoritmo es una asignación $X$, y opcionalmente el valor de
cada criterio (§\ref{sec:objetivo}) para auditoría.

\subsection{Restricciones duras}
\begin{definicion}[Factibilidad]
$X$ es \emph{factible} para $\mathcal{I}$ si cumple, para todo $e\in E$:
\begin{restriccion}[Elegibilidad]\label{r:eleg}
$g\in X_e \Rightarrow c(g)\in A(e)$.
\end{restriccion}
\begin{restriccion}[A lo sumo un grupo por materia]\label{r:uno}
$\card{X_e\cap G_c}\le 1$ para toda $c\in C$.
\end{restriccion}
\begin{restriccion}[Sin choques de horario]\label{r:choque}
Para toda $s\in S$: $\card{\{g\in X_e : s\in T(g)\}}\le 1$.
\end{restriccion}
\begin{restriccion}[Disponibilidad]\label{r:disp}
$g\in X_e\Rightarrow T(g)\cap U(e)=\emptyset$.
\end{restriccion}
\begin{restriccion}[Carga de créditos]\label{r:cred}
$\kappa_{\min}(e)\le\sum_{g\in X_e}\mathrm{cr}(c(g))\le\kappa_{\max}(e)$.
\end{restriccion}
y, para todo $g\in G$:
\begin{restriccion}[Cupo]\label{r:cupo}
$\card{X^g}\le q(g)$.
\end{restriccion}
\end{definicion}
Se denota $\mathcal{F}(\mathcal{I})$ el conjunto de asignaciones factibles.
Puede ser vacío (p.\,ej. si $\kappa_{\min}(e)$ es inalcanzable); un
algoritmo debe poder reportarlo.

\begin{obs}
Verificar la factibilidad de una $X$ dada es polinomial en el tamaño de la
instancia, y es independiente de cómo se obtuvo $X$. Este verificador sirve
como oráculo de pruebas (ver estrategia de datos).
\end{obs}

%--------------------------------------------------------------------
\section{Criterios de calidad y función objetivo}\label{sec:objetivo}
%--------------------------------------------------------------------
Dada una asignación $X$ y un estudiante $e$, sea
$O_e^\delta=\{s\in\bigcup_{g\in X_e}T(g) : \mathrm{dia}(s)=\delta\}$
el conjunto de franjas que $e$ ocupa el día $\delta$, y
$O_e=\bigcup_\delta O_e^\delta$. Por la restricción~\ref{r:choque}, cada
$s\in O_e$ pertenece a un único grupo de $X_e$, con salón
$r_e(s)$ bien definido. Sea $s_1<\dots<s_m$ el orden por periodo de $O_e^\delta$.

\paragraph{C1. Prioridad (cobertura de materias necesarias) --- se maximiza.}
\[
\mathrm{Prio}(X)=\sum_{e\in E}\ \sum_{g\in X_e}\rho(e,c(g)),
\qquad N_1=\sum_{e\in E}\sum_{c\in A(e)}\rho(e,c).
\]

\paragraph{C2. Social --- se maximiza.}
\[
\mathrm{Soc}(X)=\sum_{\{e,e'\}\subseteq E}\omega(e,e')\,\card{X_e\cap X_{e'}},
\qquad N_2=\sum_{\{e,e'\}\subseteq E}\omega(e,e')\,\card{A(e)\cap A(e')}.
\]
(Cuenta los grupos que dos estudiantes con afinidad comparten.)

\paragraph{C3. Profesor --- se maximiza.}
\[
\mathrm{Prof}(X)=\sum_{e\in E}\ \sum_{g\in X_e}\pi(e,p(g)),
\qquad N_3=\sum_{e\in E}\card{A(e)}.
\]

\paragraph{C4. Logística (desplazamientos) --- se minimiza.}
Para dos franjas ocupadas consecutivas $s_i,s_{i+1}$ de un mismo día, el
tiempo disponible para trasladarse es
$\beta+L\,(\mathrm{per}(s_{i+1})-\mathrm{per}(s_i)-1)$ minutos. Se penaliza el
déficit:
\[
\mathrm{Log}(X)=\sum_{e\in E}\ \sum_{\delta\in D}\ \sum_{i=1}^{m-1}
\max\Bigl\{0,\ d\bigl(r_e(s_i),r_e(s_{i+1})\bigr)-\beta-L\bigl(\mathrm{per}(s_{i+1})-\mathrm{per}(s_i)-1\bigr)\Bigr\},
\]
con $m=\card{O_e^\delta}$ y $N_4=d_{\max}\,\card{E}\,n_d\,(n_h-1)$, donde
$d_{\max}=\max_{r,r'}d(r,r')$.

\paragraph{C5. Tiempos libres --- se minimiza.}
Un \emph{hueco} es una franja libre entre dos ocupadas del mismo día.
\[
\mathrm{Free}(X)=\sum_{e\in E}\Bigl[\ \sum_{\delta\in D}\sum_{i=1}^{m-1}\bigl(\mathrm{per}(s_{i+1})-\mathrm{per}(s_i)-1\bigr)\ +\ \card{O_e\cap F(e)}\ \Bigr],
\]
\[
N_5=\sum_{e\in E}\bigl(n_d\max\{n_h-2,0\}+\card{F(e)}\bigr).
\]
El primer sumando penaliza huecos; el segundo penaliza ocupar franjas que el
estudiante prefería libres.

\begin{definicion}[Función objetivo]
\[
Q(X)=\lambda_1\frac{\mathrm{Prio}(X)}{N_1}+\lambda_2\frac{\mathrm{Soc}(X)}{N_2}
+\lambda_3\frac{\mathrm{Prof}(X)}{N_3}
-\lambda_4\frac{\mathrm{Log}(X)}{N_4}-\lambda_5\frac{\mathrm{Free}(X)}{N_5}.
\]
Cada cociente está en $[0,1]$ (C3 en $[-1,1]$), de modo que los pesos
$\lambda_i$ son comparables entre sí.
\end{definicion}

%--------------------------------------------------------------------
\section{Enunciado del problema}
%--------------------------------------------------------------------
\begin{definicion}[Problema de optimización: \textsc{AsigHor}]
\textbf{Entrada:} una instancia $\mathcal{I}$. \\
\textbf{Salida:} $X^\ast\in\mathcal{F}(\mathcal{I})$ tal que
$Q(X^\ast)=\max_{X\in\mathcal{F}(\mathcal{I})}Q(X)$, o el mensaje
``infactible'' si $\mathcal{F}(\mathcal{I})=\emptyset$.
\end{definicion}

\begin{definicion}[Versión de decisión: \textsc{AsigHor-D}]
Dados $\mathcal{I}$ y $\theta\in\R$, ¿existe $X\in\mathcal{F}(\mathcal{I})$ con
$Q(X)\ge\theta$?
\end{definicion}

\begin{obs}[Dificultad --- por demostrar]
Con un solo estudiante, $\lambda_1=1$ y sin criterios acoplados, el problema
se reduce a elegir a lo sumo un grupo por materia sin choques maximizando
prioridad, lo que parece corresponder a \emph{Job Interval Selection} (cada
materia es un trabajo y cada grupo una alternativa de intervalos), conocido
como NP-difícil. Se debe escribir la reducción formal; se espera que
\textsc{AsigHor-D} sea NP-completo. Esto justificaría el uso de
heurísticas o metaheurísticas con garantías empíricas en la entrega final.
\end{obs}

%--------------------------------------------------------------------
\section{Ejemplo pequeño}
%--------------------------------------------------------------------
$D=\{1,2\}$, $H=\{1,2,3,4\}$, $L=60$, $\beta=10$.
$E=\{e_1,e_2\}$, $C=\{c_1,c_2\}$ con $A(e_i)=C$, $\mathrm{cr}\equiv 3$,
$\kappa_{\min}=0$, $\kappa_{\max}=6$, $U=F=\emptyset$. Salones $r_1,r_2,r_3$
con $d(r_2,r_3)=8$ min.

\begin{center}
\begin{tabular}{lllll}
\toprule
Grupo & Materia & Profesor & Cupo & Sesiones $(\mathrm{día},\mathrm{periodo})$ en salón \\
\midrule
$g_1$ & $c_1$ & $p_1$ & 1 & $(1,1),(1,2)$ en $r_1$ \\
$g_2$ & $c_1$ & $p_2$ & 2 & $(1,3),(1,4)$ en $r_2$ \\
$g_3$ & $c_2$ & $p_1$ & 2 & $(1,2),(2,2)$ en $r_3$ \\
\bottomrule
\end{tabular}
\end{center}

\begin{itemize}[nosep]
  \item $X=\{(e_1,g_2),(e_1,g_3),(e_2,g_2),(e_2,g_3)\}$ es factible.
        Para $e_1$: franjas $(1,2)\to(1,3)$ consecutivas, déficit
        $\max\{0,\,8-10-0\}=0$.
  \item $\{(e_1,g_1),(e_1,g_3)\}\subseteq X$ viola~\ref{r:choque}
        (ambos ocupan $(1,2)$).
  \item $\{(e_1,g_1),(e_2,g_1)\}\subseteq X$ viola~\ref{r:cupo} ($q(g_1)=1$).
  \item $\{(e_1,g_1),(e_1,g_2)\}\subseteq X$ viola~\ref{r:uno} (dos grupos de $c_1$).
\end{itemize}

%--------------------------------------------------------------------
\section{Formato de entrada y salida (propuesta)}
%--------------------------------------------------------------------
Para que cualquier implementación sea comparable, se propone intercambiar
instancias y soluciones en JSON:
\begin{itemize}[nosep]
  \item \textbf{Entrada:} un objeto con las claves
        \texttt{estudiantes}, \texttt{materias}, \texttt{grupos},
        \texttt{profesores}, \texttt{salones}, \texttt{distancias},
        \texttt{parametros} (con $n_d,n_h,L,\beta,\lambda$), y por estudiante
        \texttt{elegibles}, \texttt{prioridad}, \texttt{no\_disponible},
        \texttt{preferir\_libre}, \texttt{creditos\_min/max},
        \texttt{afinidad\_profesor}; y la lista de aristas
        \texttt{afinidad\_social}.
  \item \textbf{Salida:} \texttt{asignacion}: lista de pares
        (estudiante, grupo); \texttt{valor\_Q}; \texttt{valores\_por\_criterio};
        \texttt{estado} $\in$ \{\texttt{factible}, \texttt{infactible}\}.
\end{itemize}

%--------------------------------------------------------------------
\section{Decisiones abiertas del grupo}
%--------------------------------------------------------------------
\begin{enumerate}[label=(D\arabic*)]
  \item \pend{¿Las materias obligatorias deben ser restricción dura (cubrir
        todas las de prioridad máxima) en lugar de sólo pesar en C1? Hacerlas
        duras puede volver infactible la instancia.}
  \item \pend{Granularidad de franjas: ¿1 h o bloques de 2 h como en la
        oferta real? Define $n_h$ y $L$.}
  \item \pend{¿Criterio social también cuenta amigos en la misma materia
        aunque en grupos distintos (afinidad parcial)?}
  \item \pend{Valores por defecto de $\lambda$ y calibración con
        encuestas o con conjuntos de prueba.}
  \item \pend{¿Prerrequisitos y correquisitos entre materias se modelan en
        $A(e)$ (actual) o como restricciones explícitas?}
  \item \pend{Validar con el profesor el alcance: ¿una instancia es un
        semestre completo de la universidad o una facultad/subconjunto?}
\end{enumerate}

\end{document}
